\chapter{超荷、闭合与超空间}

\lectureoverview{超对称之所以困难，不在于其规则繁多，而在于这些规则把普通物理学中彼此分离的观念联系在了一起。我们从格拉斯曼变量开始，其奇元素反交换，其函数在有限多项后终止。它们的导数代数随后为奇对称生成元提供了模型。一个玻色振子和一个费米振子给出最简单的超荷：它们的反交换子就是哈密顿量，因此交换玻色子与费米子与时间平移不可分割。同一代数可以通过一个称为超空间的扩大时空中的格拉斯曼方向上的平移来几何地表示。}

\section{为什么在前进之前必须先退后}

超对称是物理学中最抽象的构造之一。它类似于旋转和平移，因为它是一种对称性，但如果不扩大我们对坐标的含义，它就无法容纳于普通时空之中。它的自然参数是格拉斯曼量：这些对象可以相加和相乘，但反交换而非交换。

这就是为什么直接攻击最终的代数毫无帮助。我们将首先使奇算术变得可用，然后回顾普通连续对称性是如何组织的，只有到那时才把这两个观念结合起来。这条路线更长，但每一步都为下一步做准备。

\section{实用形式的格拉斯曼代数}

\subsection{奇元素与偶元素}

设\(\epsilon_i\)表示独立的格拉斯曼生成元。它们满足
\begin{equation}
\epsilon_i\epsilon_j=-\epsilon_j\epsilon_i,
\qquad
\{\epsilon_i,\epsilon_j\}=0.
\end{equation}
令\(i=j\)，得到
\begin{equation}
2\epsilon_i^2=0,
\qquad
\epsilon_i^2=0,
\end{equation}
在实数或复数上成立。生成元不为零；只有它的平方为零。

每个齐次元素都有一个宇称。包含奇数个格拉斯曼生成元的乘积是奇的，而包含偶数个的乘积是偶的：
\begin{align}
\text{odd}\times\text{odd}&=\text{even},\\
\text{odd}\times\text{even}&=\text{odd},\\
\text{even}\times\text{even}&=\text{even}.
\end{align}
偶元素与齐次元素交换。两个奇元素反交换。因此当\(n\)为偶时，
\begin{equation}
n\epsilon_i=\epsilon_i n,
\qquad
\epsilon_i\epsilon_j=-\epsilon_j\epsilon_i,
\end{equation}。

偶元素不必是普通的实数或复数。例如乘积\(\epsilon_1\epsilon_2\)是偶的且幂零。普通数构成代数的标量部分，而偶乘积提供更多交换元素。

\begin{classroomqa}
\audiencequestion{格拉斯曼代数与格拉斯曼数之间有区别吗，还是它们是同一回事？\lecturetimestamp{00:06:45}}
\lecturerresponse{元素就是格拉斯曼数；代数则是这些元素连同其加法与乘法规则的集合。因此可以宽松地说它们是同一对象，但区分被操作的事物与操作它们的规则是有用的。}
\end{classroomqa}

\subsection{函数终止}

一个格拉斯曼变量的函数只有两种可能的幂：
\begin{equation}
f(\epsilon)=a+b\epsilon.
\end{equation}
不存在\(\epsilon^2\)项。对于两个独立生成元，
\begin{equation}
f(\epsilon_1,\epsilon_2)
=a+b_1\epsilon_1+b_2\epsilon_2+c\epsilon_1\epsilon_2.
\end{equation}
展开在两个生成元都出现后停止。对于\(n\)个独立生成元，有序单项式构成一个\(2^n\)维基：
\begin{equation}
1,\quad
\epsilon_i,\quad
\epsilon_i\epsilon_j,\quad
\ldots,\quad
\epsilon_1\epsilon_2\cdots\epsilon_n.
\end{equation}
\editorialnote{有限的\(2^n\)维陈述是单生成元和双生成元例子背后的一般模式。也可以定义无限格拉斯曼代数，但这里使用的每个多项式都只涉及有限多个生成元。}

当两个实生成元组合成一个复对\(\theta,\bar\theta\)时，一般函数可以按选定的排序写成
\begin{equation}
f(\theta,\bar\theta)
=a+\bar b\,\theta+\bar\theta\,b+c\,\bar\theta\theta.
\end{equation}
横线表示这里使用的复对记号。颠倒最后的顺序会改变一个符号：
\begin{equation}
\bar\theta\theta=-\theta\bar\theta.
\end{equation}

通常方便的做法是给整个函数指定一个确定的宇称。如果\(f\)是偶的，那么\(a\)和\(c\)是偶的，而\(b\)和\(\bar b\)是奇的。如果\(f\)是奇的，这些指定就反过来。系数本身可以属于一个更大的格拉斯曼代数，只要它们不依赖于我们据以展开的那些特定变量。

\begin{classroomqa}
\audiencequestion{格拉斯曼变量可以代入三角函数或其他解析函数吗？\lecturetimestamp{00:11:06}}
\lecturerresponse{可以，但幂级数会立即终止。对于一个奇生成元，
\[
\sin\epsilon=\epsilon,\qquad
\cos\epsilon=1,\qquad
e^\epsilon=1+\epsilon.
\]
熟悉的名称保留了下来，但幂次太少，函数无法发展出它们通常的解析复杂性。}
\end{classroomqa}

\section{微分就是符号记账}

我们使用左格拉斯曼导数。其定义性的局部规则是
\begin{equation}
\partial_\theta a=0,
\qquad
\partial_\theta(\theta B)=B,
\end{equation}
当\(a\)和\(B\)独立于\(\theta\)时成立。在微分之前，把相关的\(\theta\)移到左边。每次经过另一个奇因子都会改变符号。

例如，如果\(b\)是奇的，
\begin{equation}
b\theta=-\theta b,
\qquad
\partial_\theta(b\theta)=-b.
\end{equation}
导数并没有发生什么神秘的事情。负号出现在微分之前，是在重新排列奇因子时产生的。

\subsection{双变量计算}

回到
\begin{equation}
f(\theta,\bar\theta)
=a+\bar b\,\theta+\bar\theta\,b+c\,\bar\theta\theta.
\end{equation}
首先假设\(f\)是偶的。那么\(\bar b\)是奇的，而\(c\)是偶的。重新排列含有\(\theta\)的项，
\begin{align}
\bar b\,\theta&=-\theta\bar b,\\
c\,\bar\theta\theta&=-\theta c\,\bar\theta.
\end{align}
因此
\begin{equation}
\partial_\theta f=-\bar b-c\bar\theta.
\end{equation}

如果\(f\)是奇的，那么\(\bar b\)是偶的，而\(c\)是奇的。第一项不跨越任何奇系数，而最后一项\(\theta\)跨越两个奇因子\(c\)和\(\bar\theta\)。两个符号相互抵消：
\begin{equation}
\partial_\theta f=\bar b+c\bar\theta.
\end{equation}

\begin{classroomqa}
\audiencequestion{为什么在偶的例子中两个存留下来的项都带有负号？\lecturetimestamp{00:16:49}}
\lecturerresponse{第一个符号来自将\(\theta\)移过奇系数\(\bar b\)。在最后一项中，\(c\)是偶的，但\(\theta\)必须越过奇变量\(\bar\theta\)。每次重新排列都跨越一个奇因子，因此各产生一个负号。}
\end{classroomqa}

答案取决于排序约定，但数学本身并不依赖于此。如果我们从\(c\theta\bar\theta\)而不是\(c\bar\theta\theta\)开始，所显示的系数会相差一个符号，而导数公式也会相应地改变。

\begin{classroomqa}
\audiencequestion{最后一项能否写成\(\theta\)和\(\bar\theta\)互换的形式，由此产生的符号是否有更深层的意义？\lecturetimestamp{00:18:50}}
\lecturerresponse{它可以按任意一种顺序写出，但交换这两个奇变量会改变符号。人们可以将该符号吸收到系数\(c\)中。在这个阶段，其效果只是记账：顺序重要，宇称不变，而符号的改变绝不会把奇对象变成偶对象。其物理意义出现在同样的记账方式强制实现费米子统计和超对称的时候。}
\end{classroomqa}

\begin{classroomqa}
\audiencequestion{系数\(a,b,c\)只是取自复数域吗？\lecturetimestamp{00:20:47}}
\lecturerresponse{当它们所需的宇称为偶时，它们可以是普通的复数，但它们也可能包含除所选\(\theta,\bar\theta\)之外的生成元。所关注的代数通常包含所有复标量以及一组选定的奇生成元及其乘积。}
\end{classroomqa}

\begin{classroomqa}
\audiencequestion{偶函数和奇函数的导数看起来互为相反数。这是一种普遍的对称性吗？\lecturetimestamp{00:22:31}}
\lecturerresponse{并非本身如此。符号是由每个系数的宇称以及所跨越的奇因子数目机械地决定的。有时两个答案恰好相反，但可靠的规则更简单：每交换两个奇对象就贡献一个负号。}
\end{classroomqa}

\subsection{作为算符的坐标与导数}

有一个富有启发性的普通类比。让乘以\(x\)和微分\(x\)作用于任意函数\(g(x)\)。那么
\begin{align}
\partial_x\!\left(xg\right)&=g+xg',\\
x\partial_xg&=xg',
\end{align}
所以
\begin{equation}
[\partial_x,x]g=g.
\end{equation}
由于这对每个\(g\)都成立，
\begin{equation}
[\partial_x,x]=1.
\end{equation}
取\(p=-i\partial_x\)，这就是熟悉的位置-动量代数，采用\(\hbar=1\)的单位。

格拉斯曼类比使用反对易子。取
\begin{equation}
f(\theta)=a+\theta b,
\end{equation}
其中\(a\)是偶的，\(b\)是奇的，因此\(f\)是偶的。首先，
\begin{equation}
\theta\partial_\theta f=\theta b.
\end{equation}
按相反的顺序，
\begin{align}
\partial_\theta(\theta f)
&=\partial_\theta(\theta a+\theta^2b)\\
&=a.
\end{align}
将两种顺序相加得到原函数：
\begin{equation}
\left(\partial_\theta\theta+\theta\partial_\theta\right)f
=a+\theta b=f.
\end{equation}
因此
\begin{equation}
\boxed{\{\partial_\theta,\theta\}=1.}
\end{equation}
导数\(\partial_\theta\)本身是奇的。这就是为什么它与奇坐标的正则关系是反对易子。

\begin{classroomqa}
\audiencequestion{能否使用一个同时包含偶部分和奇部分的函数，而不是赋予确定的宇称？\lecturetimestamp{00:31:05}}
\lecturerresponse{这样的非齐次元素可以构造出来，但计算通常被组织成齐次的偶对象和奇对象。确定的宇称使分次乘积中的每个符号都明确无误；混合对象则通过将其分解为偶部分和奇部分来处理。}
\end{classroomqa}

\begin{classroomqa}
\audiencequestion{在普通对易子的计算中，测试函数\(g(x)\)应该是奇的吗？\lecturetimestamp{00:31:45}}
\lecturerresponse{不是。那是普通可交换坐标的普通函数。引入它只是为了展示算符恒等式\([\partial_x,x]=1\)是如何证明的，然后再用格拉斯曼对\(\partial_\theta,\theta\)重复这一论证。}
\end{classroomqa}

\begin{classroomqa}
\audiencequestion{独立的格拉斯曼生成元的数目有限制吗？\lecturetimestamp{00:32:16}}
\lecturerresponse{没有代数上的上限。对于\(n\)个生成元，基有\(2^n\)个单项式，因此记账量迅速增长，但构造本身对任何有限的\(n\)都成立。}
\end{classroomqa}

\subsection{积分}

对一个格拉斯曼变量的积分由下式确定
\begin{equation}
\int d\theta\,1=0,
\qquad
\int d\theta\,\theta=1.
\end{equation}
因此，对于\(f(\theta)=a+\theta b\)，
\begin{equation}
\int d\theta\,f(\theta)=b=\partial_\theta f(\theta).
\end{equation}
于是积分与微分执行相同的分量提取。这就是贝列辛积分。其归一化是与平移不变性和格拉斯曼分部积分相容的唯一简单选择。\editorialnote{分量规则精确地表达了讲义中关于格拉斯曼积分与微分是同一运算的陈述。}

\section{普通对称性作为基准}

考虑一个具有厄米生成元\(G_i\)的连续幺正群。在恒等元附近，
\begin{equation}
U(\alpha)=\mathbf 1+i\alpha^iG_i+O(\alpha^2),
\end{equation}
其中\(\alpha^i\)是普通的对易参数。无穷小乘法法则由李代数编码：
\begin{equation}
[G_i,G_j]=i f_{ij}{}^{k}G_k.
\end{equation}
常数\(f_{ij}{}^k\)是结构常数。封闭性意味着任意两个生成元的对易子仍是代数中已有生成元的线性组合。

对于转动，
\begin{equation}
[L_i,L_j]=i\varepsilon_{ijk}L_k.
\end{equation}
三个厄米分量可以组合成非厄米的升降算符
\begin{equation}
L_\pm=L_x\pm iL_y,
\end{equation}
它们将磁量子数\(m\)升高或降低一。复组合是有用的，同时不改变底层物理生成元的厄米性。

在一般基下，\(f_{ij}{}^k\)只有前两个指标必须是反对称的。对于具有不变正交归一度规的紧致代数，可以降低最后一个指标并选择约定使\(f_{ijk}\)完全反对称，正如转动的情形那样。\editorialnote{这一限定将普遍反对称的对易子与熟悉的转动基中可用的更强性质区分开来。}

要使\(G_i\)生成一个不含时系统的实际对称性，需有
\begin{equation}
[H,G_i]=0.
\end{equation}
等价地，\(G_i\)将每个能量本征空间映射到自身。在没有显式时间依赖时，海森堡方程给出
\begin{equation}
\frac{dG_i}{dt}=i[H,G_i]=0,
\end{equation}
因此该生成元是一个守恒荷。

普通可观测量是偶的。费米子场是奇的，但孤立的奇场本身不是可测量的可观测量；物理算符包含偶数个费米子因子的总数。超对称性改变的是生成元的宇称，而不是可观测结果所要求的宇称。

适用于所有情形的简洁记号是分级对易子
\begin{equation}
[X,Y]_{\mathrm{gr}}
=XY-(-1)^{|X||Y|}YX.
\end{equation}
若任一算符为偶，它就是普通对易子；若两者均为奇，它就是反对易子。

\section{最小的超对称振子}

\subsection{一个玻色模和一个费米模}

取单个玻色模，其产生和湮灭算符为\(a_+\)和\(a_-\)：
\begin{equation}
[a_-,a_+]=1,
\qquad
[a_+,a_+]=[a_-,a_-]=0.
\end{equation}
其数算符为
\begin{equation}
N_B=a_+a_-,
\end{equation}
其占据数可以是\(0,1,2,\ldots\)。

对于单个费米模，引入\(c_+\)和\(c_-\)：
\begin{equation}
\{c_-,c_+\}=1,
\qquad
\{c_+,c_+\}=\{c_-,c_-\}=0.
\end{equation}
因此
\begin{equation}
c_+^2=c_-^2=0,
\qquad
N_F=c_+c_-,
\end{equation}
占据数只能是\(0\)或\(1\)。玻色算符是偶的，费米算符是奇的，且每个\(a_\pm\)与每个\(c_\pm\)对易。

\editorialnote{幂零性将费米模限制为一个占据态。玻色模的占据不受限制；课堂讨论中一个简短的相反说法由振子代数本身解决。}

假设玻色子和费米子具有相同的静止能量\(m\)，并定义
\begin{equation}
Q=\sqrt m\,a_+c_-,
\qquad
\bar Q=Q^\dagger=\sqrt m\,a_-c_+.
\end{equation}
算符\(Q\)湮灭一个费米子并产生一个玻色子。其共轭湮灭一个玻色子并产生一个费米子。每个都包含一个费米算符，因此是奇的。

在数态上，它们的作用是显式的：
\begin{align}
Q\ket{n,1}
&=\sqrt{m(n+1)}\,\ket{n+1,0},
&
Q\ket{n,0}&=0,\\
\bar Q\ket{n,0}
&=\sqrt{mn}\,\ket{n-1,1},
&
\bar Q\ket{n,1}&=0.
\end{align}
\editorialnote{数态公式是讲义中所述粒子替换作用的直接归一化形式。}

\subsection{反对易子在能量上封闭}

寻找李代数的分级类比是计算
\begin{align}
\{Q,\bar Q\}
&=Q\bar Q+\bar Q Q\\
&=m\left(a_+c_-a_-c_+
        +a_-c_+a_+c_-\right).
\end{align}
由于\(a\)是偶的，它们穿过\(c\)时不带符号：
\begin{equation}
\{Q,\bar Q\}
=m\left(a_+a_-c_-c_+
        +a_-a_+c_+c_-\right).
\end{equation}
现在使用
\begin{equation}
c_-c_+=1-c_+c_-,
\qquad
a_-a_+=a_+a_-+1.
\end{equation}
于是
\begin{align}
\{Q,\bar Q\}
&=m\left[
N_B(1-N_F)+(N_B+1)N_F
\right]\\
&=m(N_B+N_F).
\end{align}
在实时重排过程中出现的临时额外因子二消失了：一个贡献提供\(N_B\)，另一个提供\(N_F\)。

对于等质量的静止粒子，哈密顿量为
\begin{equation}
H=m(N_B+N_F).
\end{equation}
因此
\begin{equation}
\boxed{\{Q,\bar Q\}=H.}
\end{equation}

\begin{figure}[H]
\centering
\includegraphics[width=0.62\textwidth]{lecture_05_figure_02.png}
\lectureframe{00:55:46.300}
\caption{黑板上单独列出的中心封闭关系。花括号和\(Q,\bar Q,H\)是可见的；逗号由标准反对易子记号提供。}
\end{figure}

能量的可加零点已被选定，使真空态具有\(H=0\)。在超对称振子中这是自然的约定：玻色和费米零点贡献相消。\editorialnote{讲义直接使用总静止能量，并未显示单独的零点项；所述约定使这一选择明确。}

哈密顿量是偶的。它同时也是一个对称性生成元：它生成时间平移。因此，除非把偶的时间平移生成元也包括进来，否则奇生成元之间无法自行闭合。

\begin{classroomqa}
\audiencequestion{该关系式中哪一项是偶的？是哈密顿量吗？\lecturetimestamp{00:58:24}}
\lecturerresponse{是的。能量是普通的可观测量，所以\(H\)是偶的。超荷\(Q\)是奇的，而两个奇量的乘积是偶的；因此，\(\{Q,\bar Q\}\)可以等于\(H\)。一个奇生成元与一个偶生成元之间由对易子联系，而非反对易子。}
\end{classroomqa}

\subsection{完成代数}

同宇称的关系式是直接的：
\begin{equation}
\{Q,Q\}=2Q^2=0,
\qquad
\{\bar Q,\bar Q\}=2\bar Q^2=0,
\end{equation}
因为\(Q^2\)包含\(c_-^2\)，而\(\bar Q^2\)包含\(c_+^2\)。

剩下的是与\(H\)的关系式。令
\begin{equation}
H\ket{E}=E\ket{E}.
\end{equation}
如果玻色子与费米子质量相等，\(Q\)将其中一个替换为另一个而不改变能量。因此\(Q\ket{E}\)要么为零，要么是另一个具有相同本征值的本征态：
\begin{equation}
H Q\ket{E}=E Q\ket{E}.
\end{equation}
另一方面，
\begin{equation}
QH\ket{E}=Q(E\ket{E})=E Q\ket{E}.
\end{equation}
相减得，
\begin{equation}
[H,Q]\ket{E}=0.
\end{equation}
如果能量本征态构成一个完备基，则此式作为算符恒等式成立：
\begin{equation}
[H,Q]=0,
\qquad
[H,\bar Q]=0.
\end{equation}

\begin{classroomqa}
\audiencequestion{如果玻色子与费米子质量不相等，会发生什么？\lecturetimestamp{01:01:16}}
\lecturerresponse{那么将其中一个替换为另一个会改变能量，因此所提出的超荷与哈密顿量不对易，简单的超对称性就丧失了。对于
\[
H=m_BN_B+m_FN_F,
\]
显式结果为
\[
[H,Q]=(m_B-m_F)Q.
\]
它恰好在\(m_B=m_F\)时为零。\editorialnote{不等质量的对易子是能量变化论证的直接算符形式，并使等质量条件变得明确。}}
\end{classroomqa}

\begin{classroomqa}
\audiencequestion{能量本征态论证能否更慢地重复一遍？\lecturetimestamp{01:01:37}}
\lecturerresponse{从\(H\ket{E}=E\ket{E}\)开始。等质量意味着\(Q\ket{E}\)具有相同的确定能量\(E\)，所以\(HQ\ket{E}=EQ\ket{E}\)。另一种顺序给出\(QH\ket{E}=QE\ket{E}=EQ\ket{E}\)，因为\(E\)是普通标量。它们的差在每个能量本征态上都为零。完备性随后将该陈述从这些态提升为算符方程\([H,Q]=0\)。}
\end{classroomqa}

因此完整的玩具代数就是
\begin{align}
\{Q,\bar Q\}&=H,\\
\{Q,Q\}&=\{\bar Q,\bar Q\}=0,\\
[H,Q]&=[H,\bar Q]=0.
\end{align}
超荷是守恒的奇荷。更令人惊讶的是，它们的代数将一种改变粒子类型的操作与时空平移的生成元混合在一起。

\section{为什么哈密顿量意味着时间平移}

在这个简化模型中，没有空间坐标，但有时间。薛定谔方程为
\begin{equation}
i\frac{\partial}{\partial t}\ket{\psi(t)}
=H\ket{\psi(t)}.
\end{equation}

\begin{figure}[H]
\centering
\includegraphics[width=0.56\textwidth]{lecture_05_figure_03.png}
\lectureframe{01:07:26.300}
\caption{时间导数写在\(\ket{\psi}\)旁边。手遮挡了分母的一部分；完整的薛定谔方程已在上方给出。}
\end{figure}

对于有限平移\(\delta\)，
\begin{equation}
\ket{\psi(t+\delta)}
=e^{-iH\delta}\ket{\psi(t)}.
\end{equation}
展开到一阶，
\begin{equation}
\ket{\psi(t+\delta)}
=\left(\mathbf 1-i\delta H\right)\ket{\psi(t)}
+O(\delta^2).
\end{equation}
因此\(H\)是在时间中更新一个态的无穷小生成元。

诸如同位旋和色这样的内部对称性也会改变粒子标签。同位旋可以混合质子和中子态；色旋转混合红、绿、蓝夸克。然而，这些内部变换的任何组合都不会在空间或时间中平移一个系统。超对称性则不同：
\begin{equation}
\begin{gathered}
\text{odd particle-changing transformation}\\
\mathrel{\stackrel{\{\ ,\ \}}{\Longrightarrow}}\\
\text{time translation}
\end{gathered}
\end{equation}
这就是一个线索：奇生成元本身可能就是平移，但沿着与通常坐标不同的坐标。

\section{超空间}

\subsection{时间和奇坐标的函数}

引入一个函数
\begin{equation}
\Psi=\Psi(\theta,\bar\theta,t).
\end{equation}
它可以被视为玩具模型中的波函数，或者更一般地，被视为一个超场。由于奇展开会终止，
\begin{equation}
\Psi(\theta,\bar\theta,t)
=A(t)+\bar\chi(t)\theta
+\bar\theta\chi(t)+F(t)\bar\theta\theta.
\end{equation}
这一个对象将偶和奇的分量函数打包在一起。\editorialnote{此分量展开应用了先前发展的格拉斯曼多项式；它是玩具超空间上函数的标准显式内容。}

普通的时间平移保持\(\theta,\bar\theta\)不变，并使得
\begin{equation}
t\longmapsto t+\delta.
\end{equation}
展开到一阶，
\begin{equation}
\Psi(\theta,\bar\theta,t+\delta)
=\Psi(\theta,\bar\theta,t)
+\delta\,\partial_t\Psi(\theta,\bar\theta,t).
\end{equation}

\begin{figure}[H]
\centering
\includegraphics[width=0.72\textwidth]{lecture_05_figure_04.png}
\lectureframe{01:13:15.060}
\caption{黑板将\(t\to t+\delta\)置于\(\Psi\)的一阶变化之上，并直接标出了导数项。}
\end{figure}

\subsection{沿奇方向的平移}

现在发明一个平移奇坐标的变换：
\begin{equation}
\theta\longmapsto\theta+\epsilon,
\end{equation}
其中\(\epsilon\)是另一个奇格拉斯曼参数。同时，让时间获得一个偶的幂零位移，
\begin{equation}
t\longmapsto t+i\bar\theta\epsilon.
\end{equation}
\(\bar\theta\)和\(\epsilon\)都是奇的，所以它们的乘积是偶的，可以一致地加到偶坐标\(t\)上。

将参数移到左侧后，
\begin{equation}
i\bar\theta\epsilon=-i\epsilon\bar\theta.
\end{equation}
因此诱导出的变分为
\begin{align}
\delta_\epsilon\Psi
&=\epsilon\,\partial_\theta\Psi
-i\epsilon\bar\theta\,\partial_t\Psi\\
&=\epsilon q_0\Psi,
\end{align}
其中
\begin{equation}
q_0=\partial_\theta-i\bar\theta\partial_t.
\end{equation}
这个奇微分算符是组合平移的生成元。使用左导数的一个相容共轭生成元为
\begin{equation}
\bar q_0=-\partial_{\bar\theta}
+i\theta\partial_t.
\end{equation}

确切的符号取决于使用左格拉斯曼导数还是右格拉斯曼导数，以及共轭如何反转奇因子。不能依赖于约定的是时间平移下的闭合性。根据上述定义，
\begin{align}
\{q_0,\bar q_0\}
&=i\{\partial_\theta,\theta\}\partial_t
+i\{\bar\theta,\partial_{\bar\theta}\}\partial_t\\
&=2i\partial_t.
\end{align}
导数-导数项消失，奇坐标的二次项也消失。两个混合项各贡献 \(i\partial_t\)。

由于在薛定谔绘景态上 \(H=i\partial_t\)，
\begin{equation}
\{q_0,\bar q_0\}=2H.
\end{equation}
如果我们归一化
\begin{equation}
q=\frac{q_0}{\sqrt2},
\qquad
\bar q=\frac{\bar q_0}{\sqrt2},
\end{equation}
那么
\begin{equation}
\boxed{\{q,\bar q\}=H.}
\end{equation}
\editorialnote{本讲有意处理了不确定的符号，然后注意到一个因子二。所展示的左导数约定保留了这一课堂计算，同时固定了一个自洽的归一化。其他标准约定则在坐标平移和生成元之间移动符号。}

\begin{classroomqa}
\audiencequestion{超空间反对易子中最后的因子二从何而来？ \lecturetimestamp{01:24:21}}
\lecturerresponse{有两个相等的混合贡献。一个来自 \(\{\partial_\theta,\theta\}=1\)，另一个来自 \(\{\partial_{\bar\theta},\bar\theta\}=1\)。每个都乘以时间导数。因此未归一化的生成元给出 \(2i\partial_t\)；每个超荷中的因子 \(1/\sqrt2\) 恢复了先前的归一化。}
\end{classroomqa}

\begin{classroomqa}
\audiencequestion{该构造是否使用了 \(\epsilon\) 很小的假设？ \lecturetimestamp{01:24:51}}
\lecturerresponse{它使用的是幂零性，而不是普通的数值大小。由于 \(\epsilon^2=0\)，对单个 \(\epsilon\) 的每次展开都在一阶后终止。称其为无穷小或小是一种有用的类比，但截断是精确的。}
\end{classroomqa}

\editorialnote{格拉斯曼参数没有普通的正大小。它们作为无穷小参数的作用来自幂零性以及展开的精确终止。}

\subsection{几何意义}

该代数现在有两种等价的实现。在振子图像中，\(Q\) 和 \(\bar Q\) 将费米子替换为玻色子，将玻色子替换为费米子。在几何图像中，\(q\) 和 \(\bar q\) 通过奇方向平移一个函数，同时诱导出普通的时间平移。在两种图像中，
\begin{equation}
\begin{gathered}
\{\text{supercharge},\text{conjugate supercharge}\}\\
=\text{time-translation generator}.
\end{gathered}
\end{equation}

包含 \(t,\theta,\bar\theta\) 的广义坐标空间是超空间。在相对论性场论中，它还包含普通空间坐标，并且超荷在完整的时空动量 \(P_\mu\) 上闭合，而不仅仅在 \(H=P_0\) 上。弯曲这一结构将导向超引力。

一方面，格拉斯曼变量是玻色子和费米子的一种高效记账工具。另一方面，闭合关系说明了更深层的东西：粒子类型与时空几何之间的区分不再是绝对的。在奇方向上移动两次会产生沿普通方向的运动。自然界是否在实验上实现了这一结构，与这一结构在数学上自洽、经济且极其刚性这一事实是两回事。
